\subsubsection{Matrix Reprensentaties van \(\mathfrak{sl}(2)\)}

Voor we beginnen, moeten we eerst een paar zaken definiëren.

We hebben de matrices \(L_1, L_2, L_3\) die voldoen aan de commutator relaties:
\[[L_i,L_j] = i \sum_k \epsilon_{ijk}L_k\]

We definiëren de volgende operators:
\[L_\pm = L_1 \pm i L_2\]
\[L_0 = 2L_3\]

Met de commutator relaties:
\[[L_0, L_\pm] = \pm 2 L_\pm\]
\[[L_+, L_-] = L_0\]
En de eigenschap \(L_\pm^\dagger = L_\mp\).

En 
\[\vec{L}^2 = \hat{L}_x \hat{L}_x + \hat{L}_y \hat{L}_y + \hat{L}_z \hat{L}_z = \frac{1}{2}(L_+ L_- + L_- L_+) + \frac{L_0^2}{4}\]
\footnote{Ik denk dat we overal de hoedjes moeten gebruiken om aan te geven dat het operators zijn, maar ze worden vaak weggelaten voor de leesbaarheid.}
We hebben een basis staat \(|l,m \rangle\), die orthonormaal is:
\[\langle l',m' | l,m \rangle = \delta_{l,l'} \delta_{m,m'}\]

Omdat de operators \(L_3\) en \(\vec{L}^2\) commuteren, kunnen we ze gelijktijdig diagonalizeren in een specifieke basis \(\{|l,m \rangle\})\),
waarbij \(l\) en \(m\) de eigenwaarden zijn van \(\vec{L}^2\) en \(L_3\) respectievelijk:
\[
\begin{cases}
\vec{L}^2 |l,m \rangle = l(l+1) |l,m \rangle \\
L_3 |l,m \rangle = m |l,m \rangle
\end{cases}
\]

De acties van de operators \(L_3\) en \(\vec{L}^2\) op de basis staten zijn:
\[
\begin{cases}
\langle l',m'| L_3 |l,m \rangle = m \delta_{l,l'} \delta_{m,m'} \\
\langle l',m'| \vec{L}^2 |l,m \rangle = l(l+1) \delta_{l,l'} \delta_{m,m'}
\end{cases}
\]

Verder hebben we nog de eigenschappen van de ladder operators \(L_\pm\), deze anhileren de hoogste en laagste gewicht staten:
\[
\begin{cases}
L_+ |l,M \rangle = 0 \\
L_- |l,M' \rangle = 0
\end{cases}
\]
waarbij \(M\) en \(M'\) de hoogste en laagste waarden van \(m\) zijn, respectievelijk.
Als ze niet anhileren, dan verhogen of verlagen ze \(m\) met 1, en veranderen ze de waarde van \(l\) niet.
\[L_\pm |l,m \rangle = N_\pm | l , m \pm 1 \rangle\]
We hebben dat:
\begin{align*}
    |N_\mp|^2
    &= N_\mp^\dagger N_\mp 
    = \langle l,m | L_\mp^\dagger L_\mp | l,m \rangle 
    = \langle l,m | L_\pm L_\mp | l,m \rangle \\
    &= \langle l,m | \vec{L}^2 - L_3^2 \pm L_3 | l,m \rangle 
    = l(l+1) - m^2 \pm m \\
    &= l(l+1) - m(m \mp 1)
\end{align*}
Dus 
\[ N_\pm = \sqrt{l(l+1) - m(m \pm 1)} \]

Waar we gebruik hebben gemaakt van:
\begin{align*}
    L_\pm L_\mp
    &= \left(L_1 \pm i L_2\right) \left(L_1 \mp i L_2\right)\\
    &= L_1^2 + L_2^2 \pm i [L_1, L_2] \\
    &= L_1^2 + L_2^2 (L_3^2 - L_3^2) \pm L_3 \\
    &= \vec{L}^2 - L_3^2 \pm L_3
\end{align*}
De actie van de ladder operators op de basis staten is als volgt:
\[\langle l',m'| L_\pm |l,m \rangle = \sqrt{l(l+1) - m(m \pm 1)} \delta_{l,l'} \delta_{m\pm 1,m'}\]



